\chapter{Discussion}
\label{chap:discussion}

\section{Interpretation of Results}
\label{sec:discussion-interpretation}

The combination of the optimised clear path and dynamic array achieves a far greater reduction in non-strong processing time than either optimisation individually. The optimised clear path eliminates the enqueue step for queue-less weak references, routing them through a lighter processing path; once that lighter path is active, the dynamic array makes traversal of the resulting list faster through improved cache locality~\cite{drepper07-memory}. Each optimisation therefore amplifies the benefit of the other, producing a super-additive interaction. The near-equivalence of \texttt{clear\_path\_dyn} and \texttt{all} follows from this: when the clear path and dynamic array are already active, adding the separate-list mechanism provides little further gain because most of the processing work has already been removed. The \texttt{sep\_only} variant confirms this from the other direction: routing references to a separate list without changing either the traversal structure or the per-reference work yields no meaningful speedup.

The \texttt{weak\_fields} variant achieves a lower non-strong processing time than the baseline but falls short of \texttt{clear\_path\_dyn} because it bypasses the \texttt{WeakReference} pipeline entirely. The residual non-strong processing time for \texttt{weak\_fields} reflects only reference types that remain in the pipeline, such as phantom and final references, rather than any optimisation of weak-reference traversal.

Despite the substantial reductions in non-strong processing time, the \texttt{WeakReference} optimisations have a limited measurable effect on total collection time. Reference processing accounts for only a fraction of the total GC cycle; the remainder is dominated by concurrent marking, relocation, and other phases that the optimisations do not touch. The overlapping distributions for major collection and old-generation time across \texttt{WeakReference} variants reflect this: once reference processing is no longer the dominant cost, further gains in that subphase translate to at most a small fraction of total cycle time.

The bimodal shape of many violin plots for overall collection metrics suggests that the benchmark encounters two distinct operating regimes. A plausible explanation is that some GC cycles carry a heavier concurrent mark-follow workload than others, perhaps due to object lifetimes or allocation patterns interacting with the timing of the explicit \texttt{System.gc()} call. This is consistent with similar bimodality observed in the mark-follow metric. The higher variability of the single-object benchmark relative to the multi-object benchmark follows from its design: concentrating all 10 million references into a single GC cycle amplifies sensitivity to transient system effects such as cache state, NUMA placement, and thread scheduling. The multi-object benchmark spreads the reference load across five cycles, averaging out per-cycle variation.

The large advantage of \texttt{weak\_fields} on overall collection metrics stems from eliminating \texttt{WeakReference} wrapper objects entirely. Each \texttt{WeakReference} object occupies heap space for its header, referent, queue, and discovered-list pointer fields, and with 10 million such objects in the single-object benchmark this amounts to approximately 460\,MB of additional live heap. Eliminating these objects reduces GC pressure across all phases: fewer objects to mark, fewer pages to scan during relocation, and less metadata to maintain. The optimised clear path and dynamic array reduce per-reference cost within the processing pipeline, but they cannot offset the broader cost of maintaining 10 million reference objects throughout the GC cycle.

The elevated auxiliary memory of \texttt{clear\_path\_dyn} (647\,MB median in the single-object benchmark) compared to \texttt{all} (455\,MB) reflects an interaction between the optimisations. Without the separate-list mechanism, all queue-less references share a single dynamic array that retains its full allocated capacity throughout the GC cycle. With the separate-list mechanism active, references are split across two independently managed lists, bounding the peak size of each more tightly. In all cases, the additional GC-committed memory is negligible relative to the 100\,GB fixed heap. The identical heap footprints of all \texttt{WeakReference} variants confirm that the optimisations do not affect the number of wrapper objects allocated; only \texttt{weak\_fields}, which eliminates the wrapper objects themselves, reduces Java heap usage.

\section{Future Work}
\label{sec:future-work}

Can the optimisations developed for ZGC be applied to other garbage collectors in the JVM, such as the G1 garbage collector, and if so, what performance improvements can be observed?